\documentclass[11pt]{article}
\usepackage[a4paper,margin=2.5cm]{geometry}
\usepackage{amsmath,amssymb}
\usepackage{microtype}
\usepackage[colorlinks=true,linkcolor=blue,citecolor=blue,urlcolor=blue]{hyperref}
\usepackage{enumitem}
\setlist{itemsep=2pt,topsep=4pt}
\newcommand{\dd}{\mathrm{d}}

\title{Report on\\[4pt] \emph{``Supersymmetric solutions of non-relativistic 11-dimensional supergravity''}\\[2pt] \large (C.~D.~A.~Blair, IFT-UAM/CSIC-25-26)}
\date{}

\begin{document}
\maketitle
\vspace{-2em}

\paragraph{Summary of the paper.}
The paper studies bosonic backgrounds of membrane Newton--Cartan (M2NC) supergravity, the
non-relativistic theory obtained from 11-dimensional supergravity by an M2-brane decoupling
limit. Earlier work (Bergshoeff, Blair, Lahnsteiner, Rosseel 2024) found that this theory comes
in two consistent versions: a \emph{maximal} one with 32 supersymmetry parameters and a large
tower of geometric constraints, and a \emph{half-maximal} one with only 16 parameters and a
short set of constraints, but whose algebra does not close into diffeomorphisms. The author
takes M2NC limits of the M2 and M5 brane solutions in the two ways used in the recent
non-Lorentzian holography literature: (i) \emph{scaling} the number of branes by hand
($N\to c^6 n$ or $c^3 n$), and (ii) \emph{smearing} the branes along the longitudinal
directions of the limit. For each case the paper checks the constraints and equations of
motion and solves the Killing spinor equations explicitly. It also shows that any
longitudinal three-geometry with non-positive Ricci scalar solves the theory, including
AdS$_3$ and the BTZ black hole, and it treats a ``spherical'' limit of the M5 in which the
radial direction is longitudinal.

\section*{1. Principal achievement}

The main result is a clean, explicitly checked \textbf{dichotomy}:
\begin{itemize}
  \item \emph{Scaled} solutions are localised in the longitudinal directions. They satisfy the
  full constraint tower of the \emph{maximal} theory, and the equations of motion leave their
  longitudinal profile completely free.
  \item \emph{Smeared} solutions are localised in the transverse directions. They have
  non-zero $f_{Aabc}$ and $T_a{}^{\{AB\}}$, so they can \emph{only} solve the
  \emph{half-maximal} theory. Their Killing spinors carry the extra projection
  $\epsilon\to\epsilon_+$, which fits the picture of the M2NC limit itself as a
  half-BPS M2 brane in the 012 directions.
\end{itemize}
This is significant for two reasons. First, it gives the half-maximal version, which until now
looked like a formal curiosity, a physical role: if the smeared backgrounds are the right bulk
duals in the proposed non-relativistic AdS/CFT examples (as Lambert--Smith argue), then the
half-maximal theory is the one that matters for holography. This matches, independently, the
field-theory expectation that half the supersymmetries become local and drop out of the bulk.
Second, the paper tests the constraint structure of the 2024 paper on new explicit
backgrounds, which is useful because that structure rested partly on linearised evidence.

Along the way the Killing spinor analysis improves on earlier work. The paper shows that the
full (not near-horizon) scaled M2 has at least 24 Killing spinors. It shows that the
32-spinor near-horizon solution of the 2024 paper is one member of an infinite family that
exists whenever $\ln H$ is harmonic. It also identifies infinite-dimensional enhancements,
whose parameters solve a $(1+1)$-dimensional massless Dirac equation, for flat space, AdS$_3$
and the M5 limits. These enhancements plausibly match the enhanced superconformal (Virasoro-type)
symmetry seen in the dual field theories. Using dimensional reduction and T-duality, the paper
also tracks which of these Killing spinors survive in type II frames.

\section*{2. Principal drawback}

The main weakness is that the paper has \textbf{no principled notion of what counts as a
preserved supersymmetry}, and so its supersymmetry counts are not reliable statements of
physics. Several points add up to this:
\begin{itemize}
  \item \emph{The analysis depends on ansätze and is incomplete.} The Killing spinor equations
  are solved under simplifying assumptions: $\partial_0\epsilon_\pm=0$ for the M2,
  $\epsilon_-=\eta_-=0$ or $\gamma^\alpha\partial_\alpha\epsilon_+=0$ elsewhere, and
  hand-picked ``intelligent'' choices of the shift parameter $\eta_-$ guided by AdS
  spinors. Every number quoted (8, 16, 24, 32, ``infinite'') is therefore a lower
  bound. The author says so openly, but it means statements such as ``no enhancement in the
  near-horizon limit'' for the smeared M2 are less firm than the summary suggests.
  \item \emph{The emergent fermionic shift symmetries ($\eta_-,\rho_{+A}$) blur what a Killing
  spinor is.} The extra spinors (the $H^{1/3}$ modes, the $H^{5/12}$ modes, the Dirac-equation
  families) exist only because these Stückelberg-like parameters can absorb terms. Without a
  computation of the Killing superalgebra, i.e.\ spinor bilinears closing onto Killing vectors
  and checked against the bosonic isometries, one cannot tell genuine symmetries from
  artefacts of gauge redundancy. For the smeared, half-maximal solutions, which are the ones
  of holographic interest, this check is impossible in principle because the algebra does not
  close. So the claimed match with field-theory superconformal enhancement is suggestive
  rather than demonstrated.
  \item \emph{The physical status of the backgrounds is not settled.} The scaled solutions are
  under-determined: any longitudinal profile solves the equations of motion, which the author
  calls ``an embarrassment of riches''. The meaning of the by-hand rescaling of $N$ is left
  open. The spherical M5 limit is admitted to be ``quick and dirty''. The paper catalogues
  solutions well, but it cannot yet say which limit gives the physically correct bulk geometry.
\end{itemize}
There are also several typographical and minor errors that a revision should fix:
\begin{itemize}
  \item The scaled-M5 section refers to the near-horizon M5 geometry as AdS$_4\times$S$^7$;
  it should be AdS$_7\times$S$^4$.
  \item ``$\eta_+$'' should read $\eta_-$ below the AdS$_3$ Killing spinor equations.
  \item ``a differential equation for $\tilde\epsilon_+$'' should read $\tilde\epsilon_-$ in
  the scaled-M5 section.
  \item There is an index mismatch in the scaled-M5 four-form
  ($\epsilon_{IJKL}\,\dd x^I\wedge\dd x^K\wedge\dd x^L$).
  \item ``$X^m=c^{-1/2}X^m$'' should be $c^{-1/2}x^m$ in the spherical limit.
  \item The summary says the infinite M2 family appears ``only in the near-horizon limit''.
  The text actually shows it holds for any $H$ with $\ln H$ harmonic.
\end{itemize}

\section*{3. Most worthwhile follow-up: brane-sourced M2NC solutions}

Balancing feasibility and impact, I would pursue the idea the author mentions in the
conclusions: \textbf{couple the non-relativistic supergravity to a non-relativistic brane
source, and determine which limiting procedure (scaling or smearing) gives a consistently
sourced solution.}

\paragraph{Why it is high-impact.} This one calculation addresses the deepest open issue the
paper raises. It would decide whether the by-hand rescaling of $N$ can be part of a
well-defined coupled action $S_{\rm SUGRA}+N\,S_{\rm wv}$, or whether only the smeared
(transversely localised) backgrounds survive. That in turn decides which version of the
theory, maximal or half-maximal, the non-relativistic holographic proposals of Lambert--Smith,
Fontanella--Nieto García and Blair--Lahnsteiner--Obers--Yan actually live in. A source term
would also break the under-determination of the scaled solutions, because a delta-function
source fixes $H$, and so it would show whether the ``embarrassment of riches'' is physical.

\paragraph{Why it is feasible.} The ingredients already exist:
\begin{itemize}
  \item The bosonic M2NC action, and its reduction to the D2NC (``Matrix 2-brane theory'')
  limit, are known.
  \item The relevant string worldsheet action in that limit has been constructed
  (Gomis--Yan).
  \item The explicit backgrounds and their curvature data are computed in this paper.
\end{itemize}
A first step is a purely bosonic check, done in the reduced D2NC frame where the M2
configuration becomes a fundamental string, of the coupled equations of motion including the
delta-function source. This is mostly a $c$-power-counting analysis followed by a direct
calculation. Two refinements could follow:
\begin{enumerate}
  \item Lift the analysis to the M2 and M5 worldvolume actions in M2NC backgrounds, oriented
  as in this paper (that is, \emph{not} aligned with the longitudinal directions).
  \item Examine $\kappa$-symmetry of the source action. Whether $\kappa$-symmetry requires the
  maximal or only the half-maximal constraints would give an independent, physical criterion
  for which supersymmetries are preserved. That would also help with the counting problem
  described in Section~2.
\end{enumerate}

\paragraph{Risks and fallback.} The source analysis could show that neither limit is
consistently sourced without additional ingredients, such as a modified worldvolume limit or
extra background fields. That outcome would still be informative, since it would expose an
inconsistency in current non-relativistic holography proposals. A lower-risk companion
project is to compute the Killing superalgebra (spinor bilinears and their closure) for the
maximal-theory solutions, especially the M5 and AdS$_3$ cases with Dirac-equation
enhancements. This would test whether the infinite-dimensional enhancements organise into a
super-Virasoro-type algebra that matches the dual field theories.

\end{document}
